\documentclass[11pt,letterpaper]{article}
\usepackage[margin=0.86in,headheight=15pt]{geometry}
\usepackage[T1]{fontenc}
\usepackage{lmodern,amsmath,amssymb,amsthm,mathtools,bm,mathrsfs}
\usepackage{microtype,booktabs,longtable,array,enumitem,fancyhdr}
\usepackage[dvipsnames]{xcolor}
\usepackage{xurl,hyperref}
\hypersetup{colorlinks=true,linkcolor=MidnightBlue,citecolor=MidnightBlue,urlcolor=MidnightBlue,
 pdftitle={Higher Quadratic Spectral Jets},pdfauthor={Research continuation prepared for the ProveIt programme with OpenAI assistance}}
\urlstyle{tt}
\numberwithin{equation}{section}
\newtheorem{theorem}{Theorem}[section]
\newtheorem{proposition}[theorem]{Proposition}
\newtheorem{corollary}[theorem]{Corollary}
\newtheorem{lemma}[theorem]{Lemma}
\theoremstyle{definition}
\newtheorem{definition}[theorem]{Definition}
\newtheorem{remark}[theorem]{Remark}
\newcommand{\Li}{\operatorname{Li}}
\newcommand{\Res}{\operatorname{Res}}
\newcommand{\FP}{\operatorname{FP}}
\newcommand{\Reg}{\operatorname{Reg}}
\newcommand{\Rea}{\operatorname{Re}}
\newcommand{\Ima}{\operatorname{Im}}
\newcommand{\dd}{\,\mathrm{d}}
\newcommand{\C}{\mathbb C}
\newcommand{\N}{\mathbb N}
\newcommand{\Z}{\mathbb Z}
\newcommand{\file}[1]{\nolinkurl{#1}}
\newcommand{\pin}{\nolinkurl{73f18ef013d5b74eafc883efb2c15210642df12a}}
\setlength{\emergencystretch}{2em}
\setlist[enumerate]{itemsep=5pt,topsep=5pt}
\pagestyle{fancy}
\fancyhf{}
\fancyhead[L]{\small Higher quadratic spectral jets}
\fancyhead[R]{\small ProveIt research continuation}
\fancyfoot[C]{\thepage}

% BEGIN sections/00-front.tex
\title{\textbf{Higher Quadratic Spectral Jets}\\[7pt]
\large Harmonic--Stieltjes factorization defects,\\
polylogarithmic second derivatives, and exact mass collisions}
\author{A research continuation prepared for Vladimir Reshetnikov's ProveIt programme\\
\small Mathematical derivations and computational audit with OpenAI assistance}
\date{11 October 2026}
\begin{document}
\maketitle
\thispagestyle{plain}
\begin{abstract}
The first derivative of a zeta-regularized quadratic spectrum can admit an
anomaly-free factorization even when none of its higher derivatives does.
We determine the exact distinction. For every polynomial weight and every
spectral derivative order, an explicit harmonic-number coefficient formula
separates normally convergent sums from Hurwitz derivatives and generalized
Stieltjes constants. A finite local term records the collision with the
Hurwitz pole. Even weights annihilate all such local residues, but no
nonzero polynomial weight annihilates a fixed higher-order defect for every
positive mass. A joint meromorphic formula explains the discontinuous
first-derivative behavior at noninteger-weight resonances.

For the bilateral quadratic spectrum we evaluate the second derivative at
every nonpositive integer using finitely many polylogarithms, their order
derivatives, and one explicitly defined convergent polylogarithm integral
family. A Gamma normalization then gives an exact mass-antiderivative
ladder at every spectral order, with complete harmonic coefficients and
fixed boundary constants. For two distinct quadratic masses we give the
second derivative for an arbitrary polynomial weight, its fully explicit
polynomial subtraction, an ordinary one-dimensional polylogarithmic
integral, and a convergent collision expansion whose every coefficient is
a finite polygamma expression. A separate logarithmic-square identity
provides Gamma-weighted integrals and a dilogarithmic antiderivative.

These results address four concrete extension questions in the preceding
\emph{Quadratic-Spectrum Stieltjes Identities} report. The classical
multiplicative-anomaly and Bessel ingredients retain their attribution.
No minimality of the added integral coordinate, numerical period
independence, or resolution of the canonical Gaussian conjectures is
asserted. The delivery includes all-index proofs, 834 finite exact
assertions, five rejected mutation controls, 63 independent numerical
comparisons, source provenance, and additive integration notes.
\end{abstract}
\medskip
\noindent\textbf{Source anchor.} Repository metadata were read at \pin.
The matching-title predecessor TeX was read from the user's Library;
its identity with the incoming ZIP bytes was not established. No remote
repository files were changed.
\clearpage
\begingroup\small\tableofcontents\endgroup
\clearpage

% END sections/00-front.tex

% BEGIN sections/01-scope.tex
\section{Research target, conventions, and logical scope}
\label{hqj:sec:scope}

The preceding report \emph{Quadratic-Spectrum Stieltjes Identities}
\cite{predecessor} proves a first-derivative factorization formula for
polynomially weighted quadratic spectra. It classifies the weights whose
local correction vanishes identically in the mass, and obtains finite
polylogarithmic formulas at all nonpositive integer spectral points.
Its final research agenda asks for higher spectral derivatives,
higher-jet factorization defects, noninteger-weight resonances, and the
second derivative of a product with distinct masses. Those questions,
rather than another first-derivative rederivation, are the starting point
of this article.

The canonical manuscript already contains a large system of
Hurwitz--Stieltjes and Gamma identities \cite{repo}. Its source inventory
and incoming collection are too extensive for a claim that every
preserved result was audited here. The contribution below is an explicit,
proved continuation of the inspected predecessor. Some ingredients are
classical: the Hurwitz Laurent expansion, Gamma and polygamma identities,
Poisson summation for a quadratic spectrum, half-integer Bessel formulas,
and the first-derivative multiplicative-anomaly mechanism
\cite{hurwitz,gamma,polygamma,bessel,order,dowker}. We do not claim
historical priority for those ingredients or for every specialization of
the resulting identities.

\subsection{Four questions and their answers}
The predecessor's question on higher-jet defects is answered by
Theorems~\ref{hqj:thm:local} and~\ref{hqj:thm:compiler}, with a complete
vanishing classification in Theorem~\ref{hqj:thm:rigidity}. Its question
on noninteger polynomial degree receives a joint meromorphic answer and
an exact finite-part conversion in Theorem~\ref{hqj:thm:rho}.

The higher-negative-integer question is answered constructively at the
\emph{second} derivative by Theorem~\ref{hqj:thm:second}. This specifies
an adequate explicit function algebra, but does not show it is the
smallest possible one. The two-mass second-derivative and collision
question is answered for every polynomial weight in
Theorems~\ref{hqj:thm:product}--\ref{hqj:thm:collision}. Higher spectral
orders have a complete factorization-defect calculus and mass-transport
law here, not a claim of a finite Fourier reduction at every order.

\subsection{Notation and branches}
Throughout the one-sided statements, $a,b>0$ and $x_n=n+a$, $n\ge0$.
A polynomial $W(x)=\sum_{d=0}^{D}w_dx^d$ may have complex coefficients.
For bilateral statements, $0<a<1$ and the lattice is $x=n+a$, $n\in\Z$.
The positive squared masses are denoted $u=b^2$ and $v=c^2$.

All logarithms of positive numbers are real. On the right half-plane
we use the principal analytic logarithm, so for $x>0$
\[
\log(x^2+b^2)=\log(x+ib)+\log(x-ib).
\]
A logarithm of a quotient in a complex-shift identity means the
\emph{difference} of the specified logarithms; we do not silently change
branches by combining them. Parameter continuation is always from these
domains. The exponential arguments of polylogarithms in the massive
formulas lie strictly inside the unit disk.

We write $\zeta^{(r)}(s,a)=\partial_s^r\zeta(s,a)$ for order derivatives,
not derivatives in $a$, and use the standard Stieltjes convention
\begin{equation}\label{hqj:eq:stieltjes}
\zeta(1+t,a)=\frac1t+\sum_{r\ge0}\frac{(-1)^r\gamma_r(a)}{r!}t^r,
\qquad \gamma_0(a)=-\psi(a).
\end{equation}
The generalized harmonic numbers are
\[
H_m^{(r)}=\sum_{k=1}^m k^{-r},\qquad H_m=H_m^{(1)},\qquad H_0^{(r)}=0.
\]
The complete exponential Bell polynomials $Y_r$ are defined by
\begin{equation}\label{hqj:eq:bell}
\exp\left(\sum_{j\ge1}\kappa_j\frac{t^j}{j!}\right)
=\sum_{r\ge0}Y_r(\kappa_1,\ldots,\kappa_r)\frac{t^r}{r!}.
\end{equation}
A superscript $[1]$ on $\Li_p$ means $\partial_p\Li_p$, whereas
$\psi^{(j)}$ denotes the ordinary $j$th derivative of digamma.

The notation $[x^{-1}]$ extracts the coefficient in the expansion at
positive infinity. A finite part is never an unspecified summation
method: the regularizing coordinate and the subtraction are displayed.
All new numerical comparisons are floating-point diagnostics. The
proofs, not those comparisons, establish the all-index statements.

% END sections/01-scope.tex

% BEGIN sections/02-local.tex
\section{All-order factorization defects: local residues and convergent terms}
\label{hqj:sec:local}

Define initially in their absolute-convergence half-planes
\begin{align}
Z_W(s;a,b)&=\sum_{n\ge0}W(x_n)(x_n^2+b^2)^{-s},\label{hqj:eq:ZW}\\
F_{W,\pm}(s;a,b)&=\sum_{n\ge0}W(x_n)(x_n\pm ib)^{-s}.
\end{align}
Their initial conditions are respectively $\Rea s>(D+1)/2$ and
$\Rea s>D+1$. Introduce the common-scale difference
\begin{align}
\mathcal D_W(s)&=\tfrac12\bigl(F_{W,+}(2s)+F_{W,-}(2s)\bigr)-Z_W(s),
\label{hqj:eq:defect-germ}\\
\Delta_r(W;a,b)&=\mathcal D_W^{(r)}(0)
=2^{r-1}\bigl(F_{W,+}^{(r)}(0)+F_{W,-}^{(r)}(0)\bigr)-Z_W^{(r)}(0).
\label{hqj:eq:defect}
\end{align}
The factors of two are part of the definition and are essential from
second order onward.

Set
\begin{equation}\label{hqj:eq:logs}
\ell_\pm(x)=\log(1\pm ib/x),\qquad
L(x)=\log(1+b^2/x^2),\qquad
D_b(x)=\ell_+(x)-\ell_-(x)=2i\arctan(b/x).
\end{equation}
The letter $D_b$ is a logarithmic difference, not the degree of $W$.
Pointwise, the summand of \eqref{hqj:eq:defect-germ}, after removing
$W(x)x^{-2s}$, is
\begin{equation}\label{hqj:eq:cosh}
e^{-sL(x)}\bigl(\cosh(sD_b(x))-1\bigr)
=\sum_{k\ge2}h_k(x)s^k,
\end{equation}
where
\begin{equation}\label{hqj:eq:hk}
h_k(x)=\sum_{q=1}^{\lfloor k/2\rfloor}
\frac{D_b(x)^{2q}(-L(x))^{k-2q}}{(2q)!(k-2q)!}.
\end{equation}
In particular,
\[
h_2=\tfrac12D_b^2,\qquad h_3=-\tfrac12LD_b^2,
\qquad h_4=\tfrac14L^2D_b^2+\tfrac1{24}D_b^4.
\]
Only even powers of $x^{-1}$ occur in every $h_k$.

\subsection{An explicit finite-subtraction operator}
For fixed $k$, let
\begin{equation}\label{hqj:eq:Sk}
S_k(x)=\sum_{e\ge-1}a_{k,e}x^e
\end{equation}
be the finite part of the expansion of $W(x)h_k(x)$ at infinity containing
all exponents at least $-1$. Put $R_k=a_{k,-1}$, with missing coefficients
understood to be zero. Then $W(x)h_k(x)-S_k(x)=O(x^{-2})$.
Define
\begin{align}
\mathcal H_k(t)={}&\sum_{n\ge0}\bigl(W(x_n)h_k(x_n)-S_k(x_n)\bigr)x_n^{-t}
+\sum_{e\ge-1}a_{k,e}\zeta(t-e,a),\label{hqj:eq:Hk}\\
\mathcal H_{k,\mathrm{reg}}(t)={}&\mathcal H_k(t)-R_k/t.
\end{align}
The series is normally convergent for $\Rea t>-1$. In a neighborhood of
zero the second line is holomorphic. This construction fixes the
regularizing coordinate to be $x_n=n+a$.

\begin{theorem}[Complete local/nonlocal decomposition]\label{hqj:thm:local}
The functions in \eqref{hqj:eq:ZW} and the two linear-factor functions are
holomorphic at $s=0$ after meromorphic continuation. For every $r\ge1$,
\begin{equation}\label{hqj:eq:all-local}
\boxed{\quad
\frac{\Delta_r(W;a,b)}{r!}
=\frac{R_{r+1}}2+
\sum_{k=2}^r\frac{2^{r-k}}{(r-k)!}
\mathcal H_{k,\mathrm{reg}}^{(r-k)}(0).
\quad}
\end{equation}
Every regular derivative in this formula has the ordinary convergent
representation
\begin{align}
\mathcal H_{k,\mathrm{reg}}^{(j)}(0)
={}&(-1)^j\sum_{n\ge0}\bigl(W(x_n)h_k(x_n)-S_k(x_n)\bigr)\log^j x_n
\notag\\
&+\sum_{e\ne-1}a_{k,e}\zeta^{(j)}(-e,a)
+a_{k,-1}(-1)^j\gamma_j(a).
\label{hqj:eq:regular-derivative}
\end{align}
The coefficients $a_{k,e}$ are finite polynomials in $b$ and the
coefficients of $W$. No unperformed divergent sum is present on the right.
\end{theorem}
\begin{proof}
Choose a finite head so that the remaining $x_n$ exceed $2b$. On the tail,
the inverse-$x$ expansions of the logarithms and exponentials converge
uniformly on compact sets of $s$ near zero. Subtract their terms through
inverse degree $D+1$. The resulting summand is locally uniformly
$O(x_n^{-2+2|\Rea s|})$, up to a fixed polynomial in $\log x_n$ after
any fixed number of derivatives. It is therefore normally summable for
$|\Rea s|<1/4$. The subtracted terms are finite multiples of
$\zeta(2s+j-d,a+N)$.

The only possible pole at zero is the term with $j-d=1$.
For each individual spectrum its inverse-power coefficient vanishes at
least linearly in $s$, canceling that simple pole. The degree-zero
coefficient only gives Hurwitz values at nonpositive integers and is
regular. The finite head is entire in $s$. This proves holomorphy of
the individual spectra at zero.

For the difference, \eqref{hqj:eq:cosh} starts at $s^2$. The same
subtraction argument allows coefficient extraction from
$\sum_{k\ge2}s^k\mathcal H_k(2s)$. Since
\[
\mathcal H_k(2s)=\frac{R_k}{2s}
+\sum_{j\ge0}\frac{2^j\mathcal H_{k,\mathrm{reg}}^{(j)}(0)}{j!}s^j,
\]
the coefficient of $s^r$ is exactly \eqref{hqj:eq:all-local}.
Only $k\le r+1$ can contribute. The finite head and the tail can be
recombined in \eqref{hqj:eq:Hk}; its ordinary remainder converges even
when $b>a$. Finally differentiating that normally convergent remainder
and using \eqref{hqj:eq:stieltjes} gives
\eqref{hqj:eq:regular-derivative}.
\end{proof}

\subsection{The first two orders, and what parity really proves}
At the first order, Theorem~\ref{hqj:thm:local} recovers the predecessor's
local formula:
\begin{align}
\Delta_1(W;a,b)
&=\tfrac14[x^{-1}]W(x)D_b(x)^2\notag\\
&=\sum_{j\ge1}w_{2j-1}\frac{(-1)^j b^{2j}}j
\left(H_{2j}-\tfrac12H_j\right).
\label{hqj:eq:first-defect}
\end{align}
The sum is finite and independent of $a$. Its vanishing for every $b>0$
is equivalent to $W$ being even. For example
\[
\Delta_1(1)=0,\qquad \Delta_1(x)=-b^2,\qquad
\Delta_1(x^3)=\tfrac23b^4.
\]
These are cited predecessor results, also recovered here.

For the second derivative we obtain instead
\begin{equation}\label{hqj:eq:second-defect-reg}
\Delta_2(W)=\Reg_x\sum_{n\ge0}W(x_n)D_b(x_n)^2
-\tfrac12[x^{-1}]W(x)L(x)D_b(x)^2,
\end{equation}
where $\Reg_x$ is precisely \eqref{hqj:eq:Hk} with the constant Laurent
term extracted at $t=0$. Evenness of $W$ removes every $R_k$, because
an even inverse-power expansion cannot contain $x^{-1}$. It does
\emph{not} remove the regular parts in \eqref{hqj:eq:all-local}.

\begin{corollary}[An ordinary convergent second defect]\label{hqj:cor:atan}
For every $a,b>0$,
\begin{equation}\label{hqj:eq:atan-square}
\boxed{\Delta_2(1;a,b)=-4\sum_{n=0}^{\infty}
\arctan^2\!\left(\frac{b}{n+a}\right).}
\end{equation}
In particular this defect is strictly negative.
\end{corollary}
\begin{proof}
For $W=1$, $D_b^2=O(x^{-2})$, so no subtraction or residue occurs.
Insert $D_b=2i\arctan(b/x)$ in \eqref{hqj:eq:second-defect-reg}.
\end{proof}
The sign is used to give an exact counterexample, not as the principal
research target. The predecessor correctly warned against higher-jet
factorization; \eqref{hqj:eq:atan-square} makes its missing term explicit.

\begin{theorem}[Universal higher-order rigidity]\label{hqj:thm:rigidity}
Fix $a>0$ and an integer $r\ge2$. For $W\in\C[x]$,
\[
\Delta_r(W;a,b)=0\quad\hbox{for every }b>0
\quad\Longleftrightarrow\quad W=0.
\]
Thus even weights are the exact first-order kernel, but the kernel at
each fixed higher order is trivial when vanishing is required for all
positive masses.
\end{theorem}
\begin{proof}
Suppose $W\ne0$. There is an $n\ge0$ with $W(x_n)\ne0$. Continue the
finite-head/tail expression in the complex $b$-plane along a path from
the positive real axis to a punctured neighborhood of $b_0=ix_n$,
avoiding the other points $\pm ix_k$. Choose the head large enough that
the remaining tail has $x_k>|b_0|+1$. Its uniformly convergent
inverse-power representation is holomorphic near $b_0$. All other
head terms are also holomorphic there on the continued local branches.

For the $n$th term, $x_n+ib$ has a simple zero at $b_0$, while
$x_n-ib$ is nonzero. Formula \eqref{hqj:eq:defect} gives the coefficient
of the highest logarithmic power $\log^r(b-b_0)$ as
\[
(-1)^r\bigl(2^{r-1}-1\bigr)W(x_n),
\]
which is nonzero. Lower powers of that logarithm and holomorphic terms
cannot cancel it. If the original function vanished on the entire
positive real axis, its analytic continuation would vanish along the
chosen path by the identity theorem, contradicting this local
singularity. The converse for $W=0$ is immediate.
\end{proof}
This is an analytic vanishing theorem for a specified family. It is not
a theorem about algebraic independence of numerical special values.

% END sections/02-local.tex

% BEGIN sections/03-harmonic.tex
\section{A harmonic--Stieltjes coefficient formula at every derivative order}
\label{hqj:sec:harmonic}

The local/nonlocal theorem is global and coordinate-explicit. A second
form turns every derivative into a normally convergent Hurwitz series
with finite rational harmonic coefficients. It is often more convenient
for symbolic work and high-precision evaluation.

For $j\ge1$ define
\begin{equation}\label{hqj:eq:Cj}
C_j(s)=\frac{(2s)_{2j}}{(2j)!}-\frac{(s)_j}{j!}
=\sum_{k=2}^{2j}c_{j,k}s^k.
\end{equation}
The two linear terms cancel. In product form,
\begin{equation}\label{hqj:eq:C-product}
C_j(s)=\frac{s}{j}\left\{
\prod_{h=1}^{2j-1}\left(1+\frac{2s}{h}\right)
-\prod_{h=1}^{j-1}\left(1+\frac{s}{h}\right)\right\}.
\end{equation}
Consequently every $c_{j,k}$ is rational. Extend it by zero outside
$2\le k\le2j$. The first examples are
\begin{align}
C_1(s)&=2s^2,\notag\\
C_2(s)&=\tfrac23s^2(s+1)(s+2),\notag\\
C_3(s)&=\tfrac1{45}s^2(s+1)(s+2)(4s^2+18s+23),\label{hqj:eq:C-examples}\\
C_4(s)&=\tfrac1{315}s^2(s+1)(s+2)(s+3)(s+4)(2s^2+8s+11).\notag
\end{align}
The quadratic coefficient is particularly simple:
\begin{equation}\label{hqj:eq:c2}
c_{j,2}=\frac{2H_{2j-1}-H_{j-1}}j
=\frac2j\left(H_{2j}-\tfrac12H_j\right).
\end{equation}

To express all coefficients in generalized harmonic numbers, put
\[
\kappa_\ell^{(N,\alpha)}
=(-1)^{\ell-1}(\ell-1)!\alpha^\ell H_N^{(\ell)}.
\]
Expanding the logarithm of each product in \eqref{hqj:eq:C-product} gives
\begin{equation}\label{hqj:eq:c-Bell}
c_{j,k}=\frac{
Y_{k-1}(\bm\kappa^{(2j-1,2)})-
Y_{k-1}(\bm\kappa^{(j-1,1)})}{j(k-1)!},\qquad k\ge2.
\end{equation}
Only the first $k-1$ entries of each vector are used. Formula
\eqref{hqj:eq:c-Bell} is an all-order finite compiler, not a recursive
instruction to differentiate a large hypergeometric expression.

\begin{theorem}[All-order harmonic--Hurwitz--Stieltjes identity]
\label{hqj:thm:compiler}
Suppose first that $0<b<a$. Then, normally for $s$ near zero after its
removable pole cancellations,
\begin{equation}\label{hqj:eq:master-C}
\mathcal D_W(s)=\sum_{d=0}^D w_d\sum_{j\ge1}
(-1)^j b^{2j}C_j(s)\zeta(2s+2j-d,a).
\end{equation}
For every $r\ge1$,
\begin{equation}\label{hqj:eq:compiler}
\boxed{\Delta_r(W;a,b)
=r!\sum_{d=0}^D w_d\sum_{j\ge1}(-1)^j b^{2j}
\mathcal T_{j,d,r}(a),}
\end{equation}
where, with $q=2j-d$,
\begin{equation}\label{hqj:eq:Tnonres}
\mathcal T_{j,d,r}(a)=
\sum_{k=2}^r\frac{2^{r-k}c_{j,k}}{(r-k)!}
\zeta^{(r-k)}(q,a),\qquad q\ne1,
\end{equation}
and
\begin{equation}\label{hqj:eq:Tres}
\mathcal T_{j,d,r}(a)=\frac{c_{j,r+1}}2+
\sum_{k=2}^r\frac{(-2)^{r-k}c_{j,k}}{(r-k)!}
\gamma_{r-k}(a),\qquad q=1.
\end{equation}
Empty sums are zero. For arbitrary $b>0$, choose $N\ge0$ with $A=a+N>b$.
Replace $a$ by $A$ in the double series and add the finite head
\begin{align}
\sum_{n=0}^{N-1}W(x_n)\bigl\{
&2^{r-1}\bigl((-\log(x_n+ib))^r+(-\log(x_n-ib))^r\bigr)
\notag\\[-2pt]
&-(-\log(x_n^2+b^2))^r\bigr\}.
\label{hqj:eq:head}
\end{align}
This is a globally valid normally convergent representation; no
small-mass restriction remains after the head is restored.
\end{theorem}
\begin{proof}
The binomial series for the average of $(1+ib/x)^{-2s}$ and
$(1-ib/x)^{-2s}$ retains exactly the even inverse powers. Their
coefficient at $x^{-2j}$ is
$(-1)^jb^{2j}(2s)_{2j}/(2j)!$. The quadratic factor has coefficient
$(-1)^jb^{2j}(s)_j/j!$. Subtract and sum the resulting powers of $x$ in
an initial half-plane to obtain \eqref{hqj:eq:master-C}.

On a compact set of $s$, the Pochhammer ratios grow at most polynomially
in $j$; their derivatives acquire only polynomial and logarithmic
factors. For sufficiently large $j$, the Hurwitz sum and its fixed
order derivatives are bounded by a polynomial in $j$ times
$a^{-2j}$, with constants depending on the compact set and the fixed
parameters. Hence $(b/a)^{2j}$ gives normal convergence. The finitely
many exceptional Hurwitz factors are continued separately. Their only
possible pole at $s=0$ occurs at $q=1$, and $C_j(s)=O(s^2)$ cancels it.

For $q\ne1$, Taylor multiplication directly gives
\eqref{hqj:eq:Tnonres}. At $q=1$, multiply $C_j(s)$ by
\[
\zeta(1+2s,a)=\frac1{2s}
+\sum_{m\ge0}\frac{(-2)^m\gamma_m(a)}{m!}s^m.
\]
The pole raises the required coefficient index from $r$ to $r+1$,
producing \eqref{hqj:eq:Tres}. This proves the derivative formula.

When $b\ge a$, perform exactly the same argument on the tail beginning
at $A=a+N$. Its polynomial weight is still $W(x)$ in the same coordinate
$x=n+A$. The omitted ordinary finite sum is entire in $s$; taking its
$r$th derivative gives \eqref{hqj:eq:head}. Thus the global formula is
an equality of continuations, not an extrapolation of a divergent
binomial series.
\end{proof}

\subsection{Concrete consequences}
Combining \eqref{hqj:eq:atan-square}, \eqref{hqj:eq:c2}, and
\eqref{hqj:eq:compiler} gives the exact harmonic identity
\begin{equation}\label{hqj:eq:atan-harmonic}
\sum_{n\ge0}\arctan^2\!\left(\frac b{n+a}\right)
=\sum_{j\ge1}\frac{(-1)^{j+1}b^{2j}}j
\left(H_{2j}-\tfrac12H_j\right)\zeta(2j,a),\quad b<a.
\end{equation}
For $a=1$, its first terms are
\[
b^2\zeta(2)-\tfrac23b^4\zeta(4)+\tfrac{23}{45}b^6\zeta(6)
-\tfrac{44}{105}b^8\zeta(8)+\cdots.
\]
The finite-head version remains valid at every positive $b$.

For an odd weight, the Stieltjes term is visible already at the second
jet. For instance,
\begin{equation}\label{hqj:eq:xsecond}
\Delta_2(x;a,b)=4b^2\psi(a)+
4\sum_{j\ge2}\frac{(-1)^jb^{2j}}j
\left(H_{2j}-\tfrac12H_j\right)\zeta(2j-1,a),\quad b<a.
\end{equation}
At the third jet the same resonant term is
$24b^2\gamma_1(a)$; at higher orders the compiler produces
$\gamma_2(a),\gamma_3(a),\ldots$ with no additional analytic argument.
For $W=x^3$ and $r=2$, the $j=2$ term is
\[
b^4\left(2+\tfrac83\gamma_0(a)\right).
\]
The rational $2b^4$ comes from $c_{2,3}/2$ before multiplication by $2!$.
Dropping this residue contribution would give an incorrect identity.

The connection to polygamma is explicit: for integers $q\ge2$,
\begin{equation}\label{hqj:eq:polygamma-zeta}
\zeta(q,a)=\frac{(-1)^q}{(q-1)!}\psi^{(q-1)}(a).
\end{equation}
Higher spectral derivatives in \eqref{hqj:eq:Tnonres} remain ordinary
Hurwitz order derivatives. The formula does not pretend they are all
ordinary argument derivatives of polygamma.

% END sections/03-harmonic.tex

% BEGIN sections/04-resonance.tex
\section{Noninteger weights and the joint geometry of resonance}
\label{hqj:sec:rho}

For a complex parameter $\rho$, use $x^\rho=e^{\rho\log x}$ on the
positive spectrum. Let $\mathcal D(s,\rho)$ be the common-scale
linear-minus-quadratic difference with that weight. The meromorphic
continuation is most transparently written as
\begin{equation}\label{hqj:eq:rho-master}
\mathcal D(s,\rho)=\sum_{j\ge1}(-1)^jb^{2j}
C_j(s)\zeta(2s+2j-\rho,a),\qquad b<a.
\end{equation}
The same finite head as before globalizes this expression. In
particular, normal convergence is locally uniform in both $s$ and
$\rho$ away from the explicit polar divisors. The relevant positive
odd resonances are $\rho_0=2J-1$, $J\ge1$.

\begin{theorem}[Joint meromorphic normal form and finite-part conversion]
\label{hqj:thm:rho}
Near $(s,\rho)=(0,2J-1)$, with $\eta=\rho-(2J-1)$,
\begin{equation}\label{hqj:eq:rho-local}
\mathcal D(s,\rho)=
\frac{(-1)^Jb^{2J}C_J(s)}{2s-\eta}+G_J(s,\eta),
\end{equation}
where $G_J$ is jointly holomorphic and $G_J(s,\eta)=O(s^2)$.
For a nonresonant fixed $\rho$, define
$\Delta_r(\rho)=\partial_s^r\mathcal D(0,\rho)$.
Then $\Delta_1(\rho)=0$. At a resonant weight the fixed-$\rho$ derivative
instead satisfies
\begin{equation}\label{hqj:eq:rho-fp}
\boxed{\Delta_r(2J-1)
=\FP_{\rho=2J-1}\Delta_r(\rho)
+\frac{r!}{2}(-1)^Jb^{2J}c_{J,r+1}.}
\end{equation}
Here the finite part uses $\eta=\rho-(2J-1)$, and the off-resonance jets
are continued meromorphically in that variable. The extra finite term is
zero for $r\ge2J$.

Along a non-tangent ray $\eta=\kappa s$, $\kappa\ne2$, the first
coefficient is
\begin{equation}\label{hqj:eq:rho-ray}
\left.\frac{\dd}{\dd s}\mathcal D(s,2J-1+\kappa s)\right|_{s=0}
=\frac{(-1)^Jb^{2J}c_{J,2}}{2-\kappa}.
\end{equation}
The ray $\kappa=2$ lies on the polar divisor and requires subtraction
before restriction.
\end{theorem}
\begin{proof}
In \eqref{hqj:eq:rho-master}, exactly the $j=J$ term has a nearby
Hurwitz pole. Subtract $1/(2s-\eta)$ from that Hurwitz factor.
The remainder is jointly holomorphic. All other factors are locally
holomorphic, and normal convergence on a finite-head tail permits the
recombination. Every term, including the finite head, starts at $s^2$,
so $G_J=O(s^2)$. This proves \eqref{hqj:eq:rho-local} globally in the
positive mass, after the head procedure.

At $\eta\ne0$, the coefficient of $s^r$ in
$C_J(s)/(2s-\eta)$ is
\[
-\sum_{k=2}^r\frac{2^{r-k}c_{J,k}}{\eta^{r-k+1}}.
\]
It has no constant Laurent coefficient in $\eta$. At $\eta=0$, however,
the same meromorphic germ restricted before differentiation is
$C_J(s)/(2s)$, whose coefficient of $s^r$ is $c_{J,r+1}/2$.
The holomorphic remainder gives the same constant term in both orders
of operations. Multiplication by $r!(-1)^Jb^{2J}$ yields
\eqref{hqj:eq:rho-fp}. Since $\deg C_J=2J$, the last assertion about
vanishing of the local conversion follows. Restricting the normal form
to $\eta=\kappa s$ proves \eqref{hqj:eq:rho-ray}.
\end{proof}

\begin{remark}[Why a naive interpolation fails]
For nonresonant weights the first defect is zero, whereas at positive
odd weights it equals the nonzero local coefficient in
\eqref{hqj:eq:first-defect}. This is not a contradiction to analytic
continuation. Differentiation on different slices of a jointly
\emph{meromorphic} germ does not give a single jointly holomorphic
family across its intersecting divisor. At $\rho=-1$, the individual
spectra can themselves have poles at $s=0$ from their leading term;
that leading term cancels in $\mathcal D$. At that exceptional weight
we discuss derivatives of the common difference, not derivatives of
undefined individual values.
\end{remark}

The theorem resolves the resonant correction requested for $x^\rho$.
It also identifies exactly which local term is lost if the weight finite
part is taken before the spectral derivative. It does not assert
holomorphicity on the polar line or exchange two singular limits
without a subtraction convention.

% END sections/04-resonance.tex

% BEGIN sections/05-gamma.tex
\section{Logarithmic squares, Gamma integrals, and a dilogarithmic primitive}
\label{hqj:sec:gamma}

The simplest nonlocal defect has a useful independent calculus.
Let $a>0$, and let $c,d\in\C$ satisfy
$\Rea(a+c)>0$ and $\Rea(a+d)>0$. Put $h=c-d$ and
\begin{equation}\label{hqj:eq:Sdef}
\mathscr S(a;c,d)=\sum_{n\ge0}
\bigl(\log(n+a+c)-\log(n+a+d)\bigr)^2.
\end{equation}
The summand is $O(n^{-2})$, uniformly on compact sets of these
parameters, so this is an ordinary holomorphic sum. In particular
$\mathscr S(a;ib,-ib)=\Delta_2(1;a,b)$.

\begin{theorem}[A digamma logarithmic integral]\label{hqj:thm:gamma-square}
On the domain above,
\begin{equation}\label{hqj:eq:Sdigamma}
\boxed{\mathscr S(a;c,d)=2h\int_0^1
\psi(a+d+ht)\log\frac{t}{1-t}\dd t.}
\end{equation}
The integral is ordinary and convergent at both endpoints. Moreover,
\begin{equation}\label{hqj:eq:Smixed}
\partial_c\partial_d\mathscr S(a;c,d)
=-2\frac{\psi(a+d)-\psi(a+c)}{d-c},
\end{equation}
with its removable value $-2\psi'(a+c)$ when $c=d$.
Every derivative in $a$ is obtained by replacing $\psi$ by the
corresponding polygamma in \eqref{hqj:eq:Sdigamma}.
\end{theorem}
\begin{proof}
For each $n$, integrate the resolvent along the segment from $d$ to $c$:
\[
\log(n+a+c)-\log(n+a+d)
=h\int_0^1\frac{\dd t}{n+a+d+ht}.
\]
The entire segment lies in the right half-plane after adding $a$.
Compact uniform resolvent estimates justify squaring, summing, and
interchanging the two parameter integrals. The standard digamma
difference identity gives
\[
\sum_{n\ge0}\frac1{(n+A)(n+B)}
=\frac{\psi(B)-\psi(A)}{B-A}.
\]
Thus $\mathscr S$ equals
\[
h\int_0^1\int_0^1
\frac{f(v)-f(t)}{v-t}\dd v\dd t,
\qquad f(t)=\psi(a+d+ht).
\]
The divided difference has a removable diagonal and is continuous.
For a continuously differentiable $f$, symmetric exclusion of a thin
diagonal strip gives
\[
\int_0^1\int_0^1\frac{f(v)-f(t)}{v-t}\dd v\dd t
=2\int_0^1f(t)\log\frac{t}{1-t}\dd t.
\]
This follows by interchanging the two terms on the symmetric truncated
domain and then integrating $1/(t-v)$; the endpoint logarithms are
integrable and the omitted divided-difference strip tends to zero.
This proves \eqref{hqj:eq:Sdigamma}, including complex parameters.

Differentiating the normally convergent sum in $c,d$ gives
$-2\sum((n+a+c)(n+a+d))^{-1}$, which is
\eqref{hqj:eq:Smixed}. Differentiation under the integral for a fixed
number of $a$ derivatives is justified by the same right-half-plane
compactness and integrable endpoint logarithms.
\end{proof}

\subsection{An anchored primitive with explicit dilogarithms}
For the next representation assume real $c>d$ and $a+d>0$, and set $h=c-d>0$.
For $y>0$ define
\begin{equation}\label{hqj:eq:fprimitive}
f_h(y)=y\log^2(1+h/y)-2h\Li_2\!\left(\frac{h}{y+h}\right).
\end{equation}
Direct differentiation, using the ordinary dilogarithm derivative,
gives
\begin{equation}\label{hqj:eq:fderivative}
f_h'(y)=\log^2(1+h/y),\qquad
f_h(y)=-\frac{h^2}{y}+O(y^{-2})\quad(y\longrightarrow\infty).
\end{equation}
The counterterm in the next theorem is therefore necessary.

\begin{theorem}[Dilogarithmic and Gamma representations of the primitive]
\label{hqj:thm:primitive}
The series
\begin{equation}\label{hqj:eq:Pprimitive}
\mathscr P(a;c,d)=h^2\psi(a+d)+
\sum_{n\ge0}\left\{f_h(n+a+d)+\frac{h^2}{n+a+d}\right\}
\end{equation}
is normally convergent for real $a+d>0$ on compact subintervals, and
\begin{equation}\label{hqj:eq:Pderivative}
\frac{\dd}{\dd a}\mathscr P(a;c,d)=\mathscr S(a;c,d).
\end{equation}
For $a_0+d,a_1+d>0$,
\begin{align}
\mathscr P(a_1;c,d)-\mathscr P(a_0;c,d)
={}&2h\int_0^1
\log\frac{\Gamma(a_1+d+ht)}{\Gamma(a_0+d+ht)}
\log\frac{t}{1-t}\dd t.
\label{hqj:eq:gamma-dilog}
\end{align}
Thus \eqref{hqj:eq:gamma-dilog} is an exact identity between an ordinary
log-Gamma integral and an explicitly convergent dilogarithmic sum.
\end{theorem}
\begin{proof}
The asymptotic formula in \eqref{hqj:eq:fderivative} leaves an
$O(n^{-2})$ remainder in \eqref{hqj:eq:Pprimitive}, and differentiating
that remainder gives $O(n^{-3})$ on compact sets. Its derivative is
\[
h^2\psi'(a+d)+\sum_{n\ge0}
\left\{\log^2(1+h/(n+a+d))-\frac{h^2}{(n+a+d)^2}\right\}.
\]
Since $\psi'(z)=\sum_{n\ge0}(n+z)^{-2}$, the rational terms cancel
and give \eqref{hqj:eq:Pderivative}. Integrate
\eqref{hqj:eq:Sdigamma} between $a_0$ and $a_1$. Fubini is valid on the
compact $a$ interval, and the primitive of $\psi$ is the real
$\log\Gamma$. This gives \eqref{hqj:eq:gamma-dilog} with the displayed
anchoring, not an unspecified additive constant.
\end{proof}

For $|h|<a+d$, the same family has an independent harmonic expansion
\begin{equation}\label{hqj:eq:S-harmonic}
\mathscr S(a;d+h,d)=
2\sum_{r\ge2}\frac{(-1)^rH_{r-1}}r h^r\zeta(r,a+d).
\end{equation}
This follows from
$\log^2(1+z)=2\sum_{r\ge2}(-1)^rH_{r-1}z^r/r$.
Its anchored $a$ primitive uses
\[
\int_{a_0}^{a_1}\zeta(r,a+d)\dd a
=\begin{cases}
\psi(a_1+d)-\psi(a_0+d),&r=2,\\
\dfrac{\zeta(r-1,a_0+d)-\zeta(r-1,a_1+d)}{r-1},&r>2.
\end{cases}
\]
This gives another exact bridge between the Gamma integral,
harmonic coefficients, and Hurwitz values, and supplies an independent
numerical route used in the accompanying tests.

% END sections/05-gamma.tex

% BEGIN sections/06-second.tex
\section{Every negative-integer second jet of the bilateral quadratic spectrum}
\label{hqj:sec:second}

For $0<a<1$, $b>0$, define
\begin{equation}\label{hqj:eq:E}
E(s;a,b)=\sum_{n\in\Z}\bigl((n+a)^2+b^2\bigr)^{-s},\qquad \Rea s>\tfrac12.
\end{equation}
Its Poisson representation is
\begin{align}
E(s;a,b)=\frac1{\Gamma(s)}\biggl\{&\sqrt\pi\Gamma(s-\tfrac12)b^{1-2s}
\notag\\
&+4\pi^s b^{1/2-s}\sum_{n\ge1}n^{s-1/2}\cos(2\pi na)
K_{s-1/2}(2\pi bn)\biggr\}.
\label{hqj:eq:Poisson}
\end{align}
For completeness, insert
$Q^{-s}=\Gamma(s)^{-1}\int_0^\infty t^{s-1}e^{-Qt}\dd t$
in \eqref{hqj:eq:E}. Gaussian Poisson summation converts the lattice
heat kernel into
$\sqrt{\pi/t}\sum_{k\in\Z}e^{-\pi^2k^2/t}e^{2\pi ika}$.
The zero mode integrates to the first term in braces. Pair the positive
and negative modes, and use the standard integral
$\int_0^\infty t^{\nu-1}e^{-\beta t-\gamma/t}\dd t
=2(\gamma/\beta)^{\nu/2}K_\nu(2\sqrt{\beta\gamma})$
to obtain the second. Initial absolute convergence justifies these
steps, and the exponentially convergent Bessel series then continues
the result \cite{bessel}.

The expression in braces is holomorphic at every $s=-m$, $m\ge0$.
Since $1/\Gamma(s)$ has a simple zero there,
\begin{equation}\label{hqj:eq:negative-zero}
E(-m;a,b)=0.
\end{equation}
Write
\[
L_m(a,b)=E'(-m;a,b),\qquad M_m(a,b)=E''(-m;a,b).
\]
The first jet $L_m$ is the finite family in the predecessor. The second
jet differentiates the Bessel order and cannot be obtained merely by
squaring or differentiating its already specialized value.

\subsection{An explicit half-integer order derivative}
For $m\ge0$ let
\begin{equation}\label{hqj:eq:PAB}
P_m(y)=\sum_{j=0}^m\frac{(m+j)!}{j!(m-j)!2^j}y^j,
\qquad A_m(y)=(-1)^mP_m(-y).
\end{equation}
Define polynomials $B_m$ by
\begin{equation}\label{hqj:eq:Brec}
B_0(y)=0,\quad B_1(y)=2y,\quad
B_{m+1}=B_{m-1}+(2m+1)yB_m+2yP_m\quad(m\ge1).
\end{equation}
All have degree at most $m$. Let
$E_1(x)=\int_x^\infty e^{-t}\dd t/t$ for $x>0$.

\begin{lemma}[Finite polynomial form of the Bessel order derivative]
\label{hqj:lem:Bessel}
For $x>0$,
\begin{equation}\label{hqj:eq:Korder}
\left.\partial_\nu K_\nu(x)\right|_{\nu=m+1/2}
=\sqrt{\frac\pi{2x}}e^{-x}
\left\{A_m(1/x)e^{2x}E_1(2x)+B_m(1/x)\right\}.
\end{equation}
At $\nu=-m-1/2$ the derivative has the opposite sign.
\end{lemma}
\begin{proof}
The classical base is $K_{1/2}(x)=k(x)=\sqrt{\pi/(2x)}e^{-x}$.
Set $f(x)=e^{2x}E_1(2x)$, so
$f'=2f-1/x$ and $f''-2f'=1/x^2$.
The modified Bessel equation, differentiated in its order at $\nu=1/2$,
is
\[
x^2D''+xD'-(x^2+\tfrac14)D=k.
\]
Substituting $D=kf$ reduces it exactly to
$x^2(f''-2f')=1$. The real integral representation of $K_\nu$ and
Laplace localization at its minimum give
$(\partial_\nu K_\nu)/k=O(x^{-1})$ for fixed $\nu$ as $x\to\infty$;
this also follows from the differentiated standard large-$x$ expansion.
Our $f$ has the same decay. Their difference solves the homogeneous
Bessel equation: the growing solution is excluded by decay, and the
remaining multiple of $k$ is excluded by its ratio tending to zero.
Thus \eqref{hqj:eq:Korder} holds at $m=0$.

Differentiate the order recurrence
$K_{\nu+1}=K_{\nu-1}+2\nu K_\nu/x$.
The evenness $K_{-\nu}=K_\nu$ implies oddness of its order derivative.
At $\nu=1/2$, this gives the $m=1$ polynomial $B_1=2y$ and
$A_1=y-1$. Inductively the differentiated recurrence gives
\eqref{hqj:eq:Brec}. The coefficient of $e^{2x}E_1(2x)$ satisfies the
same homogeneous recurrence and equals $(-1)^mP_m(-y)$ by the finite
recurrence for $P_m$. This proves the formula at all $m$.
\end{proof}
The Bessel recurrence and order differentiability are classical
\cite{besselrec,order}; this lemma records a convenient exact
normalization for the spectral calculation. The first polynomials are
\begin{center}
\begin{tabular}{c l l l}
\toprule
$m$ & $P_m(y)$ & $A_m(y)$ & $B_m(y)$\\
\midrule
0 & $1$ & $1$ & $0$\\
1 & $1+y$ & $y-1$ & $2y$\\
2 & $1+3y+3y^2$ & $1-3y+3y^2$ & $2y+8y^2$\\
3 & $1+6y+15y^2+15y^3$ & $-1+6y-15y^2+15y^3$
& $4y+16y^2+46y^3$\\
\bottomrule
\end{tabular}
\end{center}

\subsection{One explicit polylogarithm integral family}
For $p\in\C$ set
\begin{equation}\label{hqj:eq:J}
\mathcal J_p(a,b)=\int_1^\infty
\frac{\Li_p(e^{-2\pi bw+2\pi ia})}{w+1}\dd w.
\end{equation}
This is an ordinary absolutely convergent integral, entire in $p$.
For real $a$ and $b>0$, its series form is
\begin{equation}\label{hqj:eq:Jseries}
\mathcal J_p(a,b)=\sum_{n\ge1}
\frac{e^{2\pi ian+2\pi bn}E_1(4\pi bn)}{n^p}.
\end{equation}
Indeed, expand the polylogarithm under the integral and substitute
$t=2\pi bn(w+1)$ in each exponential integral. For $p$ in a compact
set, each $p$ derivative only adds a power of $\log n$; the factor
$e^{-2\pi b n w}$ dominates this uniformly when $b$ is bounded below
by a positive constant. This proves normal convergence and all claimed
interchanges, including parameter differentiation on smaller strips in
complex $a$.

Two exact identities make the added coordinate compatible with the
existing calculus:
\begin{equation}\label{hqj:eq:J-calculus}
\partial_a\mathcal J_p=2\pi i\mathcal J_{p-1},\qquad
\partial_b\mathcal J_p=2\pi\mathcal J_{p-1}
-\frac{\Li_p(e^{-2\pi b+2\pi ia})}{b}.
\end{equation}
The first follows termwise. For the second, differentiate
\eqref{hqj:eq:Jseries} and use
$\frac{\dd}{\dd b}E_1(4\pi bn)=-e^{-4\pi bn}/b$.
Every anchored primitive in $a$ is therefore explicit; for example
\[
\int_{a_0}^{a_1}\mathcal J_p(a,b)\dd a
=\frac{\mathcal J_{p+1}(a_1,b)-\mathcal J_{p+1}(a_0,b)}{2\pi i}.
\]
The integral family is introduced to give a complete expression. No
claim that it cannot be eliminated by another identity is made.

\subsection{The complete second-jet formula}
Let $q=e^{-2\pi b+2\pi ia}$ and, for real $p$, abbreviate
\[
U_p=\Rea\Li_p(q),\qquad V_p=\Rea\Li_p^{[1]}(q),\qquad
I_p=\Rea\mathcal J_p(a,b).
\]
Write $P_m(y)=\sum_jp_{mj}y^j$, $A_m(y)=\sum_ja_{mj}y^j$,
$B_m(y)=\sum_jb_{mj}y^j$. Define
\begin{equation}\label{hqj:eq:Lzero}
L_m^{(0)}(b)=-\frac{2^{2m+1}(m!)^2\pi}{(2m+1)!}b^{2m+1},
\qquad c_m(b)=\gamma-H_m+\log(\pi/b).
\end{equation}

\begin{theorem}[All nonpositive-integer second jets]\label{hqj:thm:second}
For every $m\ge0$, $0<a<1$, and $b>0$,
\begin{align}
M_m(a,b)={}&4L_m^{(0)}(b)
\bigl(H_{2m+1}-H_m-\log(2b)\bigr)\notag\\
&+\frac{4(-1)^m m!b^m}{\pi^m}
\sum_{j=0}^m\frac1{(2\pi b)^j}
\bigl\{p_{mj}(c_m(b)U_{m+j+1}-V_{m+j+1})\notag\\[-2pt]
&\hspace{43mm}-b_{mj}U_{m+j+1}-a_{mj}I_{m+j+1}\bigr\}.
\label{hqj:eq:second-all}
\end{align}
Thus a finite set of polylogarithms, their first order derivatives, and
$\mathcal J_1,\ldots,\mathcal J_{2m+1}$ suffices for each second jet.
\end{theorem}
\begin{proof}
Write \eqref{hqj:eq:Poisson} as $E(s)=H(s)/\Gamma(s)$.
The reciprocal Gamma expansion at $s=-m+\varepsilon$ is
\[
\frac1{\Gamma(-m+\varepsilon)}
=(-1)^mm!\left\{\varepsilon+(\gamma-H_m)\varepsilon^2
+O(\varepsilon^3)\right\}.
\]
Hence
\begin{equation}\label{hqj:eq:second-H}
M_m=2(-1)^mm!\bigl(H'(-m)+(\gamma-H_m)H(-m)\bigr).
\end{equation}
For the zero mode, its logarithmic derivative is
$\psi(-m-1/2)-2\log b$. The recurrence for digamma gives
\[
\psi(-m-\tfrac12)=-\gamma-2\log2+2H_{2m+1}-H_m.
\]
Substitution in \eqref{hqj:eq:second-H} yields the first line of
\eqref{hqj:eq:second-all}.

For a nonzero Fourier mode, differentiating
$\pi^s b^{1/2-s}n^{s-1/2}$ gives $\log(\pi n/b)$.
At $s=-m$, $K_{s-1/2}=K_{m+1/2}$ has the finite polynomial $P_m$.
Its order derivative has the minus sign in Lemma~\ref{hqj:lem:Bessel}.
The resulting mode coefficient, apart from its common prefactor, is
\[
e^{-2\pi bn}n^{-m-1}
\left\{P_m((2\pi bn)^{-1})(c_m(b)+\log n)
-B_m((2\pi bn)^{-1})
-A_m((2\pi bn)^{-1})e^{4\pi bn}E_1(4\pi bn)\right\}.
\]
Now
$\sum q^n\log n/n^p=-\Li_p^{[1]}(q)$, and the last part sums by
\eqref{hqj:eq:Jseries}. Expanding the three finite polynomials gives
\eqref{hqj:eq:second-all}. The exponentially decaying modes and their
order derivatives converge normally near $s=-m$, justifying the
preceding differentiations and summations.
\end{proof}

The lowest case is already a nontrivial exact identity:
\begin{align}
M_0(a,b)={}&8\pi b\bigl(\log(2b)-1\bigr)\notag\\
&+4\left\{(\gamma+\log(\pi/b))\Rea\Li_1(q)
-\Rea\Li_1^{[1]}(q)-\Rea\mathcal J_1(a,b)\right\}.
\label{hqj:eq:M0}
\end{align}
The formula includes the entire constant Fourier mode. It is not just
an oscillatory correction to an unevaluated spectral quantity.

An independent Hurwitz representation is available from the convergent
small-mass expansion
\begin{equation}\label{hqj:eq:Ebinomial}
E(s;a,b)=\sum_{j\ge0}\frac{(-1)^j(s)_j}{j!}b^{2j}
\bigl\{\zeta(2s+2j,a)+\zeta(2s+2j,1-a)\bigr\},
\quad b<\min(a,1-a).
\end{equation}
Taylor differentiation at $s=-m$ therefore gives an identity between
\eqref{hqj:eq:second-all} and an ordinary convergent series of Hurwitz
jets. In the code this route is implemented independently of the Bessel
polynomials and is used to test $m=0,1,2,3$.

% END sections/06-second.tex

% BEGIN sections/07-mass.tex
\section{Harmonic normalization and all anchored mass antiderivatives}
\label{hqj:sec:mass}

Let $u=b^2$ and put $J_{m,r}(a,u)=E^{(r)}(-m;a,\sqrt u)$, including
$J_{m,0}=0$. In the absolute-convergence half-plane and hence by
continuation,
\begin{equation}\label{hqj:eq:E-mass}
\partial_uE(s;a,\sqrt u)=-sE(s+1;a,\sqrt u).
\end{equation}
Differentiating $r$ times at $s=-m$ gives, for $m\ge1$,
\begin{equation}\label{hqj:eq:raw-transport}
\partial_uJ_{m,r}=mJ_{m-1,r}-rJ_{m-1,r-1}.
\end{equation}
The lower-jet term cannot be dropped. At the second order it is
$-2L_{m-1}$, which explains why the unnormalized second jets are not an
Appell-type mass family.

\begin{theorem}[Gamma-normalized all-order mass transport]
\label{hqj:thm:Appell}
For $m\ge0$ define the holomorphic spectral germ
\begin{equation}\label{hqj:eq:normalized-germ}
\mathfrak A_m(\varepsilon;a,u)
=\frac{m!\Gamma(1-\varepsilon)}{\Gamma(m+1-\varepsilon)}
E(-m+\varepsilon;a,\sqrt u)
=\frac{E(-m+\varepsilon;a,\sqrt u)}
{\prod_{k=1}^m(1-\varepsilon/k)}.
\end{equation}
Then
\begin{equation}\label{hqj:eq:Appell-germ}
\partial_u\mathfrak A_m=m\mathfrak A_{m-1}\quad(m\ge1).
\end{equation}
Consequently the coefficients
$\mathfrak A_{m,r}=\partial_\varepsilon^r\mathfrak A_m(0;a,u)$ obey
$\partial_u\mathfrak A_{m,r}=m\mathfrak A_{m-1,r}$ for every $r\ge1$.
Their complete harmonic expression is
\begin{equation}\label{hqj:eq:Appell-Bell}
\mathfrak A_{m,r}
=\sum_{j=0}^{r-1}\binom rj
Y_j\bigl(H_m,1!H_m^{(2)},\ldots,(j-1)!H_m^{(j)}\bigr)
J_{m,r-j}.
\end{equation}
For every $m\ge1$ and $u\ge0$,
\begin{align}
\mathfrak A_{m,r}(a,u)
={}&\sum_{j=0}^{m-1}\binom mj u^j\mathfrak A_{m-j,r}(a,0)
\notag\\
&+m\int_0^u(u-v)^{m-1}\mathfrak A_{0,r}(a,v)\dd v.
\label{hqj:eq:all-primitives}
\end{align}
All integration constants are fixed by
\begin{equation}\label{hqj:eq:massless-jets}
J_{m,r}(a,0)=2^r\bigl\{\zeta^{(r)}(-2m,a)
+\zeta^{(r)}(-2m,1-a)\bigr\}.
\end{equation}
\end{theorem}
\begin{proof}
Use \eqref{hqj:eq:E-mass} at $s=-m+\varepsilon$. Multiplication by the
Gamma factor in \eqref{hqj:eq:normalized-germ} cancels $m-\varepsilon$:
\[
\frac{m-\varepsilon}{\prod_{k=1}^m(1-\varepsilon/k)}
=\frac{m}{\prod_{k=1}^{m-1}(1-\varepsilon/k)}.
\]
This proves \eqref{hqj:eq:Appell-germ} as an identity of holomorphic
germs, and coefficient differentiation preserves it.

The multiplier has the convergent germ expansion
\[
\prod_{k=1}^m(1-\varepsilon/k)^{-1}
=\exp\left(\sum_{\ell\ge1}\frac{H_m^{(\ell)}}\ell
\varepsilon^\ell\right).
\]
The Bell definition \eqref{hqj:eq:bell}, followed by the product rule,
gives \eqref{hqj:eq:Appell-Bell}; its last $j=r$ term vanishes because
$J_{m,0}=0$.

The lattice gap $\min(a,1-a)>0$ makes the continued spectrum holomorphic
in $u$ near zero. One may verify this directly by the normally
convergent binomial expansion \eqref{hqj:eq:Ebinomial}, after any
needed finite spectral subtractions. At $u=0$, the spectrum is
$\zeta(2s,a)+\zeta(2s,1-a)$, proving
\eqref{hqj:eq:massless-jets}. Repeated integration of the Appell
relation from zero, with these initial values, gives
\eqref{hqj:eq:all-primitives}; the integral coefficient is
$m!/(m-1)!=m$. Continuation along $u\ge0$ encounters no zero of any
quadratic eigenvalue, so the identity holds for every positive mass.
\end{proof}

The first three normalized orders are
\begin{align}
\mathfrak A_{m,1}&=L_m,\notag\\
\mathfrak A_{m,2}&=M_m+2H_mL_m,\label{hqj:eq:normalized-low}\\
\mathfrak A_{m,3}&=J_{m,3}+3H_mM_m
+3\bigl(H_m^2+H_m^{(2)}\bigr)L_m.\notag
\end{align}
Here and below $L_m,M_m$ are evaluated at $(a,\sqrt u)$ when used as
functions of $u$. The first-jet finite expression needed with
\eqref{hqj:eq:second-all} is
\begin{equation}\label{hqj:eq:Lm}
L_m(a,b)=L_m^{(0)}(b)+2(-1)^m m!
\sum_{j=0}^m\frac{(m+j)!b^{m-j}}{j!(m-j)!4^j\pi^{m+j}}
\Rea\Li_{m+j+1}(e^{-2\pi b+2\pi ia}).
\end{equation}
It follows directly from \eqref{hqj:eq:Poisson} and the finite
half-integer formula, and agrees with the predecessor.

\subsection{Explicit second-order primitives}
Since $\mathfrak A_{0,2}=M_0$, the $r=2$ case gives
\begin{align}
&m\int_0^{b^2}(b^2-v)^{m-1}M_0(a,\sqrt v)\dd v\notag\\
&\quad=M_m(a,b)+2H_mL_m(a,b)
-\sum_{j=0}^{m-1}\binom mj b^{2j}
\bigl(M_{m-j}(a,0)+2H_{m-j}L_{m-j}(a,0)\bigr).
\label{hqj:eq:Mprimitive}
\end{align}
Every term on the right now has an explicit expression from
Theorem~\ref{hqj:thm:second} and \eqref{hqj:eq:Lm}, together with the
fixed Hurwitz boundary jets. Thus \eqref{hqj:eq:Mprimitive} evaluates
an infinite family of logarithmic polylogarithm integrals without
leaving its integration constants undetermined.

The massless constants can also be written in polylogarithmic order
coordinates. Let
\[
\mathcal C_p(a)=\tfrac12\bigl(\Li_p(e^{2\pi ia})+
\Li_p(e^{-2\pi ia})\bigr),
\]
continued in $p$ with the radial convention. For real $p$ this is the
real part, but the symmetric expression preserves holomorphy for
complex $p$. The reflected Hurwitz formula \cite{hurwitz} gives
\begin{equation}\label{hqj:eq:massless-functional}
E(s;a,0)=\frac{4\Gamma(1-2s)}{(2\pi)^{1-2s}}
\sin(\pi s)\mathcal C_{1-2s}(a)
\end{equation}
near the nonpositive integers. With
$k_m=2(-1)^m(2m)!/(2\pi)^{2m}$, differentiation yields
\begin{align}
L_m(a,0)&=k_m\mathcal C_{2m+1}(a),\notag\\
M_m(a,0)&=4k_m\left\{
(\log(2\pi)+\gamma-H_{2m})\mathcal C_{2m+1}(a)
-\mathcal C_{2m+1}^{[1]}(a)\right\}.
\label{hqj:eq:massless-second}
\end{align}
At $m=0$, the first line is $-\log(4\sin^2\pi a)$.
No nonexistent ordinary value of $\Li_1(1)$ is used, since $0<a<1$.

% END sections/07-mass.tex

% BEGIN sections/08-collision.tex
\section{Two quadratic masses: the exact second derivative and collision calculus}
\label{hqj:sec:collision}

Let $P(y)=\sum_{d=0}^Dp_dy^d\in\C[y]$ and $u,v>0$. Define
\begin{align}
\mathcal Z_P(s;a;u,v)
&=\sum_{n\in\Z}P((n+a)^2)
\bigl(((n+a)^2+u)((n+a)^2+v)\bigr)^{-s},
\label{hqj:eq:product-spectrum}\\
E_{P,u}(s;a)
&=\sum_{n\in\Z}P((n+a)^2)((n+a)^2+u)^{-s}.
\label{hqj:eq:weighted-E}
\end{align}
Their initial absolute-convergence conditions are
$\Rea s>(2D+1)/4$ and $\Rea s>D+1/2$, respectively. We use their
meromorphic continuations below. The finite polynomial identity
\begin{equation}\label{hqj:eq:polynomial-translate}
E_{P,u}(s;a)=\sum_{\ell=0}^D
\frac{P^{(\ell)}(-u)}{\ell!}E(s-\ell;a,\sqrt u)
\end{equation}
will reduce the individual weighted jets to the already evaluated family.

\subsection{The complete subtraction, not a formal divergent logarithm sum}
For $r\ge2$, put
\begin{equation}\label{hqj:eq:Tr}
T_r(u,v)=(-1)^r\sum_{k=1}^{r-1}
\frac{(u^k-v^k)(u^{r-k}-v^{r-k})}{k(r-k)}.
\end{equation}
These are the coefficients of $y^{-r}$ in
$\log^2((y+u)/(y+v))$ at infinity. The first three are
\begin{align}
T_2&=(u-v)^2,\notag\\
T_3&=-(u-v)^2(u+v),\label{hqj:eq:Tlow}\\
T_4&=(u-v)^2\left\{\tfrac23(u^2+uv+v^2)+\tfrac14(u+v)^2\right\}.
\notag
\end{align}
The full polynomial part of the weighted logarithmic square is
\begin{equation}\label{hqj:eq:PiP}
\Pi_P(y;u,v)=\sum_{d=2}^Dp_d\sum_{r=2}^dT_r(u,v)y^{d-r}.
\end{equation}
Define the ordinary sum
\begin{equation}\label{hqj:eq:KP}
\mathcal K_P(a;u,v)=\sum_{n\in\Z}\left\{
P(x^2)\log^2\frac{x^2+u}{x^2+v}-\Pi_P(x^2;u,v)\right\},
\qquad x=n+a.
\end{equation}
Its summand is $O(x^{-2})$, locally uniformly in positive $u,v$.
For degree at most one no subtraction is needed. At larger degree the
subtraction is generally nonzero and must not be discarded.

\begin{theorem}[Second derivative of a two-mass spectrum]
\label{hqj:thm:product}
Both spectra in \eqref{hqj:eq:product-spectrum}--\eqref{hqj:eq:weighted-E}
are holomorphic at $s=0$. For every polynomial $P$,
\begin{align}
\mathcal Z_P''(0;a;u,v)
={}&2\sum_{\ell=0}^D\frac1{\ell!}
\bigl\{P^{(\ell)}(-u)M_\ell(a,\sqrt u)
+P^{(\ell)}(-v)M_\ell(a,\sqrt v)\bigr\}
\notag\\
&-\mathcal K_P(a;u,v).
\label{hqj:eq:product-second}
\end{align}
When $u=v$, the correction vanishes and the formula reduces to
$\mathcal Z_P''(0;a;u,u)=4E_{P,u}''(0;a)$, as required by
$\mathcal Z_P(s;a;u,u)=E_{P,u}(2s;a)$.
\end{theorem}
\begin{proof}
On either positive lattice tail, all mass expansions and the weight
$P(x^2)$ contain only even powers of $x$. A pole of the associated
Hurwitz factor at $s=0$ would require a power $x^{-1}$, which is
impossible. Finite-head subtractions therefore prove holomorphy at
zero exactly as in Section~\ref{hqj:sec:local}, without any residue
term.

Compare $\mathcal Z_P(s)$ with
$\tfrac12(E_{P,u}(2s)+E_{P,v}(2s))$. After taking out $x^{-4s}$,
the pointwise difference has zero constant and linear coefficients,
and its quadratic coefficient is
\[
\tfrac12P(x^2)\log^2\frac{x^2+u}{x^2+v}.
\]
There is no $x^{-1}$ residue to lower the spectral order. Thus its
second derivative is the continued value of the logarithmic-square
sum at the fixed lattice coordinate.

The only nonsummable terms of that sum are the polynomial
\eqref{hqj:eq:PiP}. Their bilateral regularized sums vanish exactly:
for every $k\ge0$,
\begin{equation}\label{hqj:eq:even-polynomial-zero}
\zeta(-2k,a)+\zeta(-2k,1-a)=0.
\end{equation}
For example, this follows from
$\zeta(-2k,a)=-B_{2k+1}(a)/(2k+1)$ and
$B_{2k+1}(1-a)=-B_{2k+1}(a)$, including $k=0$.
Consequently the regularized value is exactly the ordinary remainder
\eqref{hqj:eq:KP}, not that remainder plus an unknown constant.
Taking the second derivative of the comparison and using
\eqref{hqj:eq:polynomial-translate} proves
\eqref{hqj:eq:product-second}.
\end{proof}

For example, at $P(y)=y^2$ the full correction is
\[
\mathcal K_{y^2}(a;u,v)=\sum_{n\in\Z}
\left\{x^4\log^2\frac{x^2+u}{x^2+v}-(u-v)^2\right\}.
\]
At $P(y)=y^3$ it is
\[
\sum_{n\in\Z}\left\{x^6\log^2\frac{x^2+u}{x^2+v}
-(u-v)^2x^2+(u-v)^2(u+v)\right\}.
\]
These illustrate why positivity of an unsubtracted logarithmic square
does not imply positivity of every polynomially regularized correction.

\subsection{An ordinary one-dimensional polylogarithmic reduction}
The positive first resolvent has the elementary formula
\begin{equation}\label{hqj:eq:E1}
E(1;a,\sqrt t)
=\frac\pi{\sqrt t}
\frac{\sinh(2\pi\sqrt t)}{\cosh(2\pi\sqrt t)-\cos(2\pi a)}.
\end{equation}
One proof takes the difference of the cotangent partial fractions at
$a\pm i\sqrt t$; another differentiates the known $L_0$ in the squared
mass. Put
\begin{equation}\label{hqj:eq:RP}
\mathcal R_P(t)=P(-t)E(1;a,\sqrt t).
\end{equation}

\begin{theorem}[The two-mass integral and mixed derivative]
\label{hqj:thm:Kintegral}
For $u>v>0$,
\begin{equation}\label{hqj:eq:Kintegral}
\boxed{\mathcal K_P(a;u,v)=2\int_v^u
\mathcal R_P(t)\log\frac{u-t}{t-v}\dd t.}
\end{equation}
The endpoint singularities are integrable logarithms. The mixed
parameter derivative is
\begin{equation}\label{hqj:eq:Kmixed}
\partial_u\partial_v\mathcal K_P(a;u,v)
=-2\frac{\mathcal R_P(u)-\mathcal R_P(v)}{v-u},
\end{equation}
with the continuous coincident value $2\mathcal R_P'(u)$.
Writing $u=b^2$, $v=c^2$, $b>c>0$, the integral is equivalently
\begin{align}
\mathcal K_P(a;b^2,c^2)=4\pi\int_c^b P(-y^2)
&\left\{1+2\Rea\Li_0(e^{-2\pi y+2\pi ia})\right\}
\notag\\[-2pt]
&\times\log\frac{b^2-y^2}{y^2-c^2}\dd y.
\label{hqj:eq:Kpolylog}
\end{align}
For $u<v$ use symmetry; for $u=v$ the correction is zero.
\end{theorem}
\begin{proof}
Polynomial division and \eqref{hqj:eq:even-polynomial-zero} show that
the regularized weighted resolvent is
\[
\Reg\sum_{n\in\Z}\frac{P(x^2)}{x^2+t}
=P(-t)E(1;a,\sqrt t)=\mathcal R_P(t).
\]
The convention here is precisely polynomial subtraction in $x^2$.
Partial fractions therefore give
\begin{equation}\label{hqj:eq:double-resolvent}
\Reg\sum_{n\in\Z}\frac{P(x^2)}{(x^2+t)(x^2+w)}
=\frac{\mathcal R_P(t)-\mathcal R_P(w)}{w-t}.
\end{equation}
Its diagonal is removable. Differentiating
\eqref{hqj:eq:KP} with respect to $u,v$ is valid after subtracting the
corresponding polynomial coefficients; the remainder and its
parameter derivatives remain $O(x^{-2})$ on compact sets. Since the
mixed derivative of the logarithmic square is
$-2/((x^2+u)(x^2+v))$, this proves \eqref{hqj:eq:Kmixed}.

Alternatively write the logarithm as
$\int_v^u(x^2+t)^{-1}\dd t$ and square it. Polynomial subtraction
commutes with these integrals coefficient by coefficient; the remaining
uniformly convergent sum may be integrated by ordinary Fubini. Thus
$\mathcal K_P$ is the double integral of
\eqref{hqj:eq:double-resolvent} over $[v,u]^2$.
The same symmetric divided-difference calculation used in
Theorem~\ref{hqj:thm:gamma-square}, with its sign reversed, gives
\eqref{hqj:eq:Kintegral}.

Finally substitute $t=y^2$ and use the absolutely convergent Poisson
kernel identity
\[
\frac{\sinh(2\pi y)}{\cosh(2\pi y)-\cos(2\pi a)}
=1+2\Rea\frac{e^{-2\pi y+2\pi ia}}{1-e^{-2\pi y+2\pi ia}}
=1+2\Rea\Li_0(e^{-2\pi y+2\pi ia}).
\]
This proves \eqref{hqj:eq:Kpolylog} and its convergence.
\end{proof}
The formulas give the full second product derivative in terms of
finite negative-integer second jets and a displayed ordinary integral.
There is no implicit subtraction of two divergent logarithmic sums.

\subsection{Every collision coefficient is a finite polygamma expression}
\begin{theorem}[Convergent collision expansion]\label{hqj:thm:collision}
Let $g=\operatorname{dist}(a,\Z)=\min(a,1-a)$. For $u>0$ and
$|h|<u+g^2$, the correction has the convergent Taylor expansion
\begin{align}
\mathcal K_P(a;u+h,u)
={}&2\sum_{r\ge2}\frac{(-1)^rH_{r-1}}r h^r
\sum_{\ell=0}^{\min(D,r-1)}
\frac{P^{(\ell)}(-u)}{\ell!}
E(r-\ell;a,\sqrt u).
\label{hqj:eq:collision-series}
\end{align}
It defines the analytic branch continued from $h=0$ in that disk.
In particular,
\begin{equation}\label{hqj:eq:confluent}
\lim_{h\to0}\frac{\mathcal K_P(a;u+h,u)}{h^2}
=P(-u)E(2;a,\sqrt u)+P'(-u)E(1;a,\sqrt u)
=-\mathcal R_P'(u).
\end{equation}
All resolvents on the right of \eqref{hqj:eq:collision-series} have
positive integer order and admit finite polygamma expressions.
\end{theorem}
\begin{proof}
The elementary identity
\[
\log^2(1+z)=2\sum_{r\ge2}\frac{(-1)^rH_{r-1}}r z^r,
\qquad |z|<1,
\]
follows either by multiplying the logarithm series or by differentiating
both sides. Put $z=h/(x^2+u)$. Its modulus is uniformly less than one
on every compact subdisk $|h|<u+g^2$. The polynomial subtraction in
\eqref{hqj:eq:KP} depends holomorphically on $h$, and its remainder is
normally summable on such compact sets. Taylor coefficients may
therefore be extracted after summation and subtraction.

At order $r$, polynomial expansion around $-u$ gives
\[
\frac{P(x^2)}{(x^2+u)^r}
=\sum_{\ell=0}^D\frac{P^{(\ell)}(-u)}{\ell!}(x^2+u)^{\ell-r}.
\]
For $\ell\ge r$ the terms are polynomials in $x^2$ and have zero
bilateral regularized sum. For $\ell<r$ they are ordinarily summable
positive-order resolvents. This gives exactly
\eqref{hqj:eq:collision-series}. The $r=2$ coefficient yields the
first equality in \eqref{hqj:eq:confluent}; differentiating
\eqref{hqj:eq:RP} and using $\partial_uE(1)=-E(2)$ gives the second.
\end{proof}

For clarity, a finite polygamma formula can be written without a
recursive mass derivative. For $N\ge1$, define
\begin{align}
A_{N,r}(b)&=(-1)^{N-r}\binom{2N-r-1}{N-r}(-2ib)^{-(2N-r)},
\quad 1\le r\le N,\notag\\
\Xi_1(z)&=-\psi(z),\qquad
\Xi_r(z)=\frac{(-1)^r}{(r-1)!}\psi^{(r-1)}(z)\quad(r\ge2).
\end{align}
Then the elementary partial fractions at the conjugate roots give
\begin{equation}\label{hqj:eq:finite-polygamma}
E(N;a,b)=2\Rea\sum_{r=1}^N A_{N,r}(b)
\bigl\{\Xi_r(a+ib)+\Xi_r(1-a+ib)\bigr\}.
\end{equation}
For $r=1$ the two conjugate harmonic tails cancel before the digamma
values are taken. The formula can also be checked from
\begin{equation}\label{hqj:eq:resolvent-derivative}
E(N;a,\sqrt u)=\frac{(-1)^{N-1}}{(N-1)!}
\partial_u^{N-1}E(1;a,\sqrt u).
\end{equation}
Thus no noninteger-order or unevaluated infinite spectral sum is
required in any collision coefficient.

For $P=1$, the expansion begins
\begin{align}
\mathcal K_1(a;u+h,u)
={}&h^2E(2)-h^3E(3)+\tfrac{11}{12}h^4E(4)
-\tfrac56h^5E(5)+\cdots,
\label{hqj:eq:collision-low}
\end{align}
all at $(a,\sqrt u)$. For higher-degree $P$, the extra translated
polynomial derivatives in \eqref{hqj:eq:collision-series} are essential;
keeping only $P(-u)E(r)$ would be incorrect.

\subsection{An exact averaged logarithmic-square identity}
An explicit elementary consequence provides both an independent check
and a useful integral evaluation.
\begin{corollary}[Average over the lattice shift]\label{hqj:cor:mean}
For $b,c>0$,
\begin{align}
\int_0^1\mathcal K_1(a;b^2,c^2)\dd a
&=\int_{-\infty}^{\infty}
\log^2\frac{x^2+b^2}{x^2+c^2}\dd x\notag\\
&=8\pi\bigl\{b\log(2b)+c\log(2c)
-(b+c)\log(b+c)\bigr\}.
\label{hqj:eq:mean-K}
\end{align}
Consequently
\begin{align}
\int_0^1\mathcal Z_1''(0;a;b^2,c^2)\dd a
=8\pi\bigl\{&b\log(2b)+c\log(2c)
+(b+c)\log(b+c)\notag\\
&-2(b+c)\bigr\}.
\label{hqj:eq:mean-product}
\end{align}
\end{corollary}
\begin{proof}
Absolute convergence unfolds the shift average of the sum into the
real-line integral. Alternatively the mean of $E(1;a,\sqrt t)$ is
$\pi/\sqrt t$. Averaging \eqref{hqj:eq:Kintegral} and substituting
$t=y^2$ gives, for $b>c$,
\[
4\pi\int_c^b\log\frac{b^2-y^2}{y^2-c^2}\dd y.
\]
Split the logarithm into its four linear factors. The integrals of
$\log(b-y)$ and $-\log(y-c)$ cancel. Integrating the remaining
$\log(b+y)-\log(y+c)$ gives twice the braces in
\eqref{hqj:eq:mean-K}. Symmetry covers $c>b$ and continuity covers
$b=c$. The mean of \eqref{hqj:eq:M0} is its zero Fourier mode
$8\pi b(\log(2b)-1)$, so \eqref{hqj:eq:product-second} proves
\eqref{hqj:eq:mean-product}.
\end{proof}

% END sections/08-collision.tex

% BEGIN sections/09-audit.tex
\section{Source audit, mathematical status, and reproducible evidence}
\label{hqj:sec:audit}

\subsection{What was actually inspected}
The canonical manuscript README, its directory structure, the incoming
directory listing, and repository tree metadata were retrieved through
the connected GitHub tools. The tree commit returned during this
reading was
\begin{center}\small\pin\end{center}
The incoming catalogue contained six archives. Their names and Git blob
identifiers are recorded in \file{PROVENANCE.json}. Their binary members
were not inspected: direct binary retrieval was unavailable in this
session. It would therefore be incorrect to claim a full audit of those
archives or a proof that none contains overlapping results.

The complete matching-title TeX of
\file{ProveIt_Quadratic_Spectrum_Stieltjes_2026-10-10.tex} was recovered
from the user's Library and materialized for programmatic inspection.
The targeted proof sections, finite first-jet formulas, and final
research questions were read. Its SHA-256 is
\begin{center}\small
\nolinkurl{4185cc4aa2299c8f417ca376b3ddc4af3fab58752d4cd93f70c2afad661f95aa}
\end{center}
The predecessor records its own earlier repository pin; that pin is not
replaced by the newer metadata anchor above. Matching the title does
not prove byte identity with the incoming ZIP. The source hash is
recorded so a future repository integration can resolve that provenance
question exactly.

The source's higher-jet warning was correct. No new error in the
inspected first-jet proof is asserted here. The work extends that proof
rather than silently replacing it. The canonical Gaussian $S_6$ and
revised $S_8$ conjectures are not settled by these spectral identities,
and their status must not be changed on the basis of this report.

\subsection{Corrections to tempting but invalid extensions}
The following are safeguards demonstrated by the present formulas,
not accusations that the predecessor made these mistakes.

\paragraph{First-jet parity is not higher-jet multiplicativity.}
Even $W$ eliminates the residue coefficients $R_k$, but not the
regular terms. Formula \eqref{hqj:eq:atan-square} disproves the
extension already at $W=1$, $r=2$; Theorem~\ref{hqj:thm:rigidity}
settles universal higher-order vanishing completely.

\paragraph{The colliding coefficient is one order higher.}
At $2j-d=1$, the Hurwitz pole turns $c_{j,r+1}s^{r+1}$ into a
contribution of order $s^r$. The term $c_{j,r+1}/2$ in
\eqref{hqj:eq:Tres} cannot be dropped. Weight finite parts and spectral
derivatives differ by exactly \eqref{hqj:eq:rho-fp}.

\paragraph{Finite heads matter beyond the first derivative.}
The ordinary head in \eqref{hqj:eq:head} vanishes for $r=1$, but
usually not for $r\ge2$. A calculation that merely replaces $a$ by
$a+N$ in the tail loses an explicit finite term. The tests compare
multiple head choices at masses both inside and outside the original
small-mass disk.

\paragraph{The polynomial subtraction is part of the two-mass answer.}
At degree at least two, the raw logarithmic-square sum diverges.
Equation \eqref{hqj:eq:PiP} gives every required term, and
\eqref{hqj:eq:even-polynomial-zero} explains why no further finite
constant is present. The second derivative of a product is not the
sum of its two second derivatives; the scale factor and the correction
are both shown in \eqref{hqj:eq:product-second}.

\paragraph{A sufficient function family is not a minimal one.}
The integral $\mathcal J_p$ makes the second-jet formula complete.
Neither a failed reduction search nor the appearance of $E_1$ proves
that $\mathcal J_p$ cannot be reduced further. That minimality question
remains separate.

\subsection{Executed finite exact checks}
The standalone script \file{verification/check_exact.py} completed
\textbf{834 exact assertions} using rational arithmetic and SymPy.
It also rejected \textbf{five deliberately corrupted identities}.
The exact suite independently compares the Pochhammer and harmonic
coefficient constructions through $j=20$, checks the Bessel and
order-differentiated differential equations through $m=10$, checks
Gamma-normalization identities through $m=16$, verifies logarithmic-square
coefficients through degree $18$, verifies polynomial resolvent
remainders through weight degree $10$, and checks the relevant
Bernoulli reflection and first-anomaly coefficients.

The report \file{data/exact_results.json} enumerates every assertion
rather than reporting an inferred count. The exact rational table
\file{data/coefficient_table.json} gives low-order $C_j$, all checked
$P_m,A_m,B_m$, and sample two-mass subtraction coefficients. These are
finite regression certificates. The mathematical range beyond the
checked indices follows from the proofs, not from extrapolation.

\subsection{Executed numerical diagnostics}
Three separately replayable suites completed \textbf{63 comparisons}
at 50 decimal working digits with scaled tolerance $10^{-34}$.
The retained maximum scaled discrepancy is below $8\cdot10^{-49}$.
The factorization suite has 19 comparisons, the Fourier/Bessel suite
24, and the collision suite 20. Their individual inputs, outputs,
and discrepancies are in the three corresponding JSON files under
\file{data/}.

The independent comparisons include separate complex-shift Hurwitz
jets versus the cancellation compiler, global digamma integrals
versus the head/tail formula, Bessel order derivatives versus finite
polynomials, ordinary polylogarithm integrals versus exponential-integral
series, and the negative-integer second jets versus the independent
quadratic binomial expansion. The collision tests compare subtracted
spectral sums with ordinary integrals and Taylor expansions. Three
additional comparisons compute the two-mass second derivative directly
from its double Pochhammer expansion, independently of the
logarithmic-square decomposition.

These are not rigorous interval enclosures. The truncation routines
choose generous precision-dependent cutoffs, and quadrature is
floating-point. The implementation avoids a numerical library overflow
at astronomically small complex polylogarithm arguments by using the
leading term of their convergent series once the discarded terms are
far below the diagnostic precision. That engineering safeguard does
not replace an error-certified numerical analysis. No proof-assistant
formalization or external peer review is claimed.

\subsection{Replay and additive integration}
From the package root, with Python, mpmath, SymPy, and a LaTeX engine:
\begin{verbatim}
python verification/verify_manifest.py
python verification/check_exact.py
python verification/check_numeric.py --part factorization \
  --output data/numerical_factorization.json
python verification/check_numeric.py --part fourier \
  --output data/numerical_fourier.json
python verification/check_numeric.py --part collision \
  --output data/numerical_collision.json
python verification/build.py
\end{verbatim}
Run manifest verification before replaying or rebuilding.
The delivered JSON files are receipts of executed runs. Replaying with
other output paths preserves those receipts. The build script assembles
a self-contained TeX source from the modular sources and runs three
serial LaTeX passes. Rebuilding can change PDF metadata; manifest
verification is for the delivered bytes, not a certification of every
future rebuild.

A suitable additive destination is
\file{Analysis/Polylogarithms/docs/reports/higher-quadratic-jets/}.
The integration notes map each theorem to the predecessor's specific
research question. All manuscript labels use the prefix \file{hqj:}.
The short integration fragment is an overview with pointers, not a
substitute for the full proofs. No remote commit, overwrite, or
automatic conjecture-status change is part of this delivery.

% END sections/09-audit.tex

% BEGIN sections/10-research.tex
\section{Further research questions and concrete next targets}
\label{hqj:sec:research}

The questions below go beyond the settings proved here. They are
formulated to keep exact evaluation and analytic validity separate from
unsupported claims of transcendental nonreduction.

\begin{enumerate}
\item \textbf{Third and higher negative-integer Fourier jets.}
For $r\ge3$, derive a terminating polynomial recurrence for the
order derivatives $\partial_\nu^{r-1}K_\nu$ at half-integers.
Determine a finite list of explicit convergent polylogarithmic
integral kernels at each fixed $r$ and give every coefficient of
$E^{(r)}(-m;a,b)$ uniformly in $m$. The present $\mathcal J_p$ is the
first additional coordinate; the normalized mass ladder supplies
consistency relations for any proposed higher-order answer.

\item \textbf{Eliminating or identifying the added integral coordinate.}
At rational $a$ or special $b$, seek exact identities for
$\mathcal J_p(a,b)$ in terms of familiar polylogarithmic, Hurwitz,
or Gamma order derivatives. Alternatively define a formal function
algebra precisely enough to prove a limited obstruction there.
Do not infer numerical nonreducibility merely from the failure of a
finite symbolic basis search.

\item \textbf{Higher-order two-mass and many-mass interactions.}
Replace the second-jet logarithmic square by the cubic and quartic
mixed logarithmic terms of a product of several quadratic factors.
Derive explicit polynomial subtractions and reduce their parameter
integrals to divided-resolvent kernels. A first useful target is the
complete third jet for two masses, with a convergent collision
expansion and correctly normalized repeated-mass limit.

\item \textbf{Several noninteger weight resonances at once.}
The local equation $2s-\eta=0$ in Theorem~\ref{hqj:thm:rho} has one
simple pole. For products with independently varied orders and
several weight parameters, classify the intersecting polar divisors
and give a multivariable finite-part conversion law. State explicitly
which transverse coordinates define the finite part and how that
choice affects the local correction.

\item \textbf{Rational-shift arithmetic projections.}
For $a=p/q$, apply finite character projectors simultaneously to the
harmonic--Stieltjes compiler, the second-jet Fourier family, and the
two-mass integral. Produce a complete small-conductor dictionary
with even and odd character components kept distinct. The goal is
an exact family of Gamma or Dirichlet-$L$ derivative identities,
not a numerical guess about the number of independent periods.

\item \textbf{Mixed shift--mass primitives.}
Combine the $a$-primitive law \eqref{hqj:eq:J-calculus} with the
normalized mass ladder to evaluate rectangular parameter integrals
of the second jets. Determine which mixed primitives close in the
same finite set of $\Li_p$, $\Li_p^{[1]}$, and $\mathcal J_p$, and
which require an additional explicitly anchored coordinate.

\item \textbf{One-sided polynomially weighted logarithmic products.}
The bilateral cancellation of even polynomial sums simplifies
Section~\ref{hqj:sec:collision}. On a one-sided lattice those
polynomial values need not vanish. Derive the full two-mass
analogue, with its Bernoulli values, digamma terms, and local
Stieltjes corrections retained. The finite subtraction framework
in Section~\ref{hqj:sec:local} already gives a precise convention
for this extension.

\item \textbf{Analytic formalization and certified evaluation.}
Formalize the common-scale pole cancellation, the finite-head
normal-convergence lemma, and the logarithmic divided-difference
integral identity in a proof assistant. Separately implement interval
arithmetic for the $\mathcal J_p$ integrals and the globally
subtracted collision sums. The existing exact coefficient tables
and mutation controls are finite fixtures, not substitutes for those
analytic developments.
\end{enumerate}

\section*{Conclusion}
\addcontentsline{toc}{section}{Conclusion}
The key distinction is between a local residue and the full higher
spectral derivative. First-order even-weight factorization survives
because the residue vanishes; higher derivatives retain a normally
convergent logarithmic contribution. The all-order compiler makes that
contribution and every Stieltjes contact term explicit. A separate
half-integer Bessel calculation completes every negative-integer second
jet, and a Gamma normalization restores exact mass integration at all
spectral orders. Finally the two-mass product has a complete ordinary
integral correction and a convergent, finite-polygamma collision
calculus. These are exact identities with stated domains and proofs,
not conjectural numerical reductions.

% END sections/10-research.tex

% BEGIN sections/11-appendix.tex
\appendix
\section{Independent Hurwitz expansions used in the verification}
\label{hqj:sec:appendix}

This appendix gives the separate analytic formulas used for the most
important numerical comparisons. They are useful identities in their
own right and make the independence of the implementations inspectable.
Let
\[
R_a(z)=\zeta(z,a)+\zeta(z,1-a),\qquad g=\min(a,1-a).
\]
All primes in this appendix mean derivatives of $R_a$ in $z$.

\subsection{The second jet at a negative integer}
For fixed $m\ge0$ and $0<b<g$, define for $0\le j\le m$
\[
h_{mj}=H_m-H_{m-j},\qquad
h_{mj}^{(2)}=H_m^{(2)}-H_{m-j}^{(2)},\qquad q_j=2j-2m.
\]
Differentiating \eqref{hqj:eq:Ebinomial} gives
\begin{align}
M_m(a,b)={}&\sum_{j=0}^m\binom mjb^{2j}
\left\{(h_{mj}^2-h_{mj}^{(2)})R_a(q_j)
-4h_{mj}R_a'(q_j)+4R_a''(q_j)\right\}
\notag\\
&+\sum_{j=m+1}^{\infty}
\frac{(-1)^{j+m}b^{2j}}{j\binom{j-1}{m}}
\left\{2(H_{j-m-1}-H_m)R_a(q_j)+4R_a'(q_j)\right\}.
\label{hqj:eq:appendix-M}
\end{align}
For $j\le m$, logarithmic differentiation of $(s)_j$ at $s=-m$
gives the finite harmonic differences. For $j>m$, separate its simple
factor $s+m$ and use
\[
\frac{(-1)^j(-m+\varepsilon)_j}{j!}
=\frac{(-1)^{j+m}}{j\binom{j-1}{m}}
\left\{\varepsilon+(H_{j-m-1}-H_m)\varepsilon^2+O(\varepsilon^3)\right\}.
\]
This proves \eqref{hqj:eq:appendix-M}. There are no colliding Hurwitz
poles at $q_j$, since every $q_j$ is even. Comparing this formula with
\eqref{hqj:eq:second-all} gives a proved Hurwitz--polylogarithm integral
identity for every $m$, independent of numerical testing.

\subsection{The two-mass second derivative without using the defect}
For $0<u,v<g^2$, expand the two quadratic factors independently:
\begin{equation}\label{hqj:eq:double-Pochhammer}
\mathcal Z_P(s;a;u,v)=
\sum_{d=0}^Dp_d\sum_{j,k\ge0}
\frac{(-1)^{j+k}(s)_j(s)_k}{j!k!}u^jv^k
R_a(4s+2j+2k-2d).
\end{equation}
Normal convergence follows from the geometric factors relative to
the lattice gap. At $s=0$, all potential Hurwitz arguments are even,
so no pole correction is required. Collecting by $\ell=j+k$ yields
\begin{align}
\mathcal Z_P''(0;a;u,v)
={}&16\sum_{d=0}^Dp_dR_a''(-2d)
\notag\\
&+\sum_{\ell\ge1}(-1)^\ell\sum_{d=0}^Dp_d
\biggl[\frac{u^\ell+v^\ell}{\ell}
\left\{2H_{\ell-1}R_a(2\ell-2d)+8R_a'(2\ell-2d)\right\}
\notag\\[-2pt]
&\hspace{28mm}+2R_a(2\ell-2d)
\sum_{j=1}^{\ell-1}\frac{u^jv^{\ell-j}}{j(\ell-j)}\biggr].
\label{hqj:eq:product-direct}
\end{align}
The $j=k=0$ term contributes the first line. Terms with exactly one
positive index use
$(s)_\ell/\ell!=s/\ell+H_{\ell-1}s^2/\ell+O(s^3)$;
terms with both indices positive use their two linear coefficients.
This proves \eqref{hqj:eq:product-direct} directly from the defining
product spectrum. The separate code route does not call
$\mathcal K_P$, the Bessel recurrence, or the product-defect theorem.
Its agreement with \eqref{hqj:eq:product-second} is therefore a useful
cross-check of signs and scale normalization.

\section{Selected symbolic and numerical receipts}
\label{hqj:sec:receipts}
The delivered receipts include all individual checks. Two sample values
from the executed factorization and second-jet diagnostics are
\begin{align*}
\sum_{n\ge0}\arctan^2\!\left(\frac{0.4}{n+0.7}\right)
&\approx0.3939890185058070693735654866107854335592,\\
\Delta_2(1;0.7,0.4)
&\approx-1.575956074023228277494261946443141734237.
\end{align*}
These displayed decimals illustrate the normalization. They are not
claimed as interval-enclosed values. The actual test files retain
50-digit working-precision values and scaled errors.

The rational constructions also detect simple implementation errors
exactly. Replacing $C_1(s)=2s^2$ by $3s^2$ fails the Pochhammer
coefficient comparison; omitting $B_1(y)=2y$ fails the differentiated
Bessel equation; omitting the lower-jet term in
\eqref{hqj:eq:raw-transport} fails the Gamma-normalization relation.
The final negative controls are intentionally false statements, not
failed mathematical claims retained as theorems.

% END sections/11-appendix.tex
\begingroup\raggedright
\begin{thebibliography}{99}
\addcontentsline{toc}{section}{References}
\bibitem{predecessor}
ProveIt research continuation,
\emph{Quadratic-Spectrum Stieltjes Identities: Complete reflected jets,
harmonic counterterms, finite polylogarithmic reductions, and a polynomial
parity criterion}, 10 October 2026. Prepared for Vladimir Reshetnikov's
ProveIt programme with OpenAI assistance. Matching-title Library TeX
read in this continuation; hash and source scope in Section~\ref{hqj:sec:audit}.
Especially its Section 5, Section 6, and Section 7.6 questions 1, 2, 3,
and 7. Incoming catalogue entry:
\url{https://github.com/VladimirReshetnikov/ProveIt/blob/main/docs/incoming/ProveIt_Quadratic_Spectrum_Stieltjes_2026-10-10.zip}.
The archive contents were not verified against the Library source.

\bibitem{repo}
ProveIt Contributors,
\emph{Polylogarithms and their Arithmetic Bridges}, canonical manuscript
README and source structure. Repository metadata anchor \pin.
\url{https://github.com/VladimirReshetnikov/ProveIt/tree/main/Analysis/Polylogarithms/docs/manuscript}.
Incoming catalogue:
\url{https://github.com/VladimirReshetnikov/ProveIt/tree/main/docs/incoming}.
Read 11 October 2026; precise scope in the provenance artifact.

\bibitem{hurwitz}
NIST Digital Library of Mathematical Functions,
\S25.11, \emph{Hurwitz Zeta Function}.
\url{https://dlmf.nist.gov/25.11}.
Used for the Laurent convention, Bernoulli values, Gamma relation,
argument derivatives, and reflected Fourier identities. Consulted
11 October 2026.

\bibitem{polylog}
NIST Digital Library of Mathematical Functions,
\S25.12, \emph{Polylogarithms}.
\url{https://dlmf.nist.gov/25.12}.
Used for defining series, analytic continuation, and derivative
conventions. Consulted 11 October 2026.

\bibitem{gamma}
NIST Digital Library of Mathematical Functions,
\S5.7, \emph{Series Expansions} for the Gamma and psi functions.
\url{https://dlmf.nist.gov/5.7}. Consulted 11 October 2026.

\bibitem{polygamma}
NIST Digital Library of Mathematical Functions,
\S5.15, \emph{Polygamma Functions}.
\url{https://dlmf.nist.gov/5.15}. Consulted 11 October 2026.

\bibitem{bessel}
NIST Digital Library of Mathematical Functions,
\S10.25, \emph{Definitions}, and \S10.32,
\emph{Integral Representations} for modified Bessel functions.
\url{https://dlmf.nist.gov/10.25}; \url{https://dlmf.nist.gov/10.32}.
Consulted 11 October 2026.

\bibitem{besselrec}
NIST Digital Library of Mathematical Functions,
\S10.29, \emph{Recurrence Relations and Derivatives}.
\url{https://dlmf.nist.gov/10.29}. Consulted 11 October 2026.

\bibitem{order}
NIST Digital Library of Mathematical Functions,
\S10.38, \emph{Derivatives with Respect to Order}.
\url{https://dlmf.nist.gov/10.38}. Consulted 11 October 2026.

\bibitem{dowker}
J.~S.~Dowker, \emph{Calculation of the multiplicative anomaly},
arXiv:1412.0549, first posted 2014, corrected version 3,
24 September 2023.
\url{https://arxiv.org/abs/1412.0549}.
Cited for the classical determinant-anomaly setting. The common-scale
higher-derivative convention and all formulas used here are derived
explicitly in the text.
\end{thebibliography}
\endgroup
\end{document}
